% Generated by scripts/assemble.py; corrections belong in source/corrections.json.
\examhead{153}{中山大学 2025 年数学分析试卷}{145}
% p146-b16; source PDF page 146
\blockanchor{p146-b16}
\question{1}填空题. 每题 5 分, 共 50 分.\par
% p146-b17; source PDF page 146
\blockanchor{p146-b17}
(1) $\lim_{n\to \infty}\sum_{k = 1}^{n}\frac{\sin\frac{k}{n}\pi}{n + \frac{k}{n}} =$\par
% p146-b18; source PDF page 146
\blockanchor{p146-b18}
(2) 设 $f(x)$ 在 $x = 0$ 处可导, 且 $f(0) = 2$, 则 $\lim_{n \to \infty} \left( \frac{f\left(\frac{2}{n}\right)}{f(0)} \right)^n =$ \_\_\_\_.\par
% p146-b19; source PDF page 146
\blockanchor{p146-b19}
(3) $\sum_{n=1}^{\infty} \frac{n(n+1)}{2^n} = \underline{\quad}$.\par
% p146-b20; source PDF page 146
\blockanchor{p146-b20}
(4) $\sum_{n=1}^{\infty} \frac{1}{3^n \sqrt{n+1} x^{3n}}$ 的收敛域为 \_\_\_\_.\par
% p146-b21; source PDF page 146
\blockanchor{p146-b21}
(5) $f(s,t) = e^{2st}\sin (s^2 +t)$ ，则 $\frac{\partial f}{\partial s} = \underline{\quad}$\par
% p146-b22; source PDF page 146
\blockanchor{p146-b22}
(6) $f(x,y,z) = x^{2} + y^{2} + z^{2}$, 则 $\operatorname{div}(\nabla f) =$ \_\_\_\_.\par
% p146-b23; source PDF page 146
\blockanchor{p146-b23}
(7) 设  $f(x)=\left\{\begin{aligned}e^{-\frac{1}{x^{2}}}, & x\neq0;\\ 0, & x=0.\end{aligned}\right.$  则  $f^{(2025)}(0)=$  \_\_\_\_.\par
% p146-b24; source PDF page 146
\blockanchor{p146-b24}
(8) $\iint_{x^2 + y^2 \leq 4} e^{-x^2 - y^2} \, \mathrm{d}x \, \mathrm{d}y = \underline{\quad}$.\par
% p147-b00; source PDF page 147
\blockanchor{p147-b00}
(9) (本题可能有误) $\oint_{L} (y - z) \mathrm{d}x + (z - x) \mathrm{d}y + (x - y) \mathrm{d}z =$ \_\_\_\_，其中 $L$ 为柱面 $x^{2} + y^{2} = R^{2}$ 与平面 $\frac{x}{R} + \frac{y}{H} = 1 (R, H > 0)$ 的交线，从 $x$ 轴看沿逆时针方向.\par
% p147-b01; source PDF page 147
\blockanchor{p147-b01}
(10) $\int_{0}^{+\infty}\frac{\sin(\sqrt{2}x)}{x}\mathrm{d}x =$\par
% p147-b02; source PDF page 147
\blockanchor{p147-b02}
\question{2}(15 分) 证明 $\{x_{n}\}$ 收敛, 其中 $x_{n} = \sum_{k=1}^{n} \frac{1}{k} - \ln n$.\par
% p147-b03; source PDF page 147
\par\Needspace{6\baselineskip}
\blockanchor{p147-b03}
\question{3}(15 分) 设 $f(x)$ 在 $[0, +\infty)$ 上可导, 且 $f'(x)$ 在 $[0, +\infty)$ 上连续, 若 $\lim_{x \to +\infty} [2f(x) + f'(x)] = 0$, 证明:\par
\nopagebreak[4]
% p147-b04; source PDF page 147
\blockanchor{p147-b04}
\displayformula{\lim _ {x \to + \infty} f (x) = 0.}
% p147-b05; source PDF page 147
\blockanchor{p147-b05}
\question{4}(15 分) 设 $f(x)$ 在 $[0,1]$ 上连续, 在 $(0,1)$ 上可导, 且 $f(1) = 2\int_{0}^{\frac{1}{2}} x^{2}f(x) \mathrm{d}x$. 证明: 存在 $\eta \in (0,1)$, 使得 $2f(\eta) + \eta f'(\eta) = 0$.\par
% p147-b06; source PDF page 147
\par\Needspace{6\baselineskip}
\blockanchor{p147-b06}
\question{5}(15 分) 设 $f(x)$ 在 $[0, 2\pi]$ 上无穷次可导, 且 $f(0) = f(2\pi)$, $\int_{0}^{2\pi} f(x) \, \mathrm{d}x = 0$. 证明:\par
\nopagebreak[4]
% p147-b07; source PDF page 147
\blockanchor{p147-b07}
\displayformula{\int_ {0} ^ {2 \pi} f ^ {2} (x) \mathrm{d} x \leq \int_ {0} ^ {2 \pi} [ f ^ {\prime} (x) ] ^ {2} \mathrm{d} x.}
% p147-b08; source PDF page 147
\blockanchor{p147-b08}
\question{6}(20 分) 设 $f(x)$ 在 $[0,1)$ 上连续, 证明: $f(x)$ 在 $[0,1)$ 上一致连续的充要条件是存在 $[0,1]$ 上的连续函数 $g(x)$, 当 $x \in [0,1)$ 时, 有 $g(x) = f(x)$.\par
% p147-b09; source PDF page 147
\blockanchor{p147-b09}
\question{7}(20 分) 解答如下问题:\par
% p147-b10; source PDF page 147
\par\Needspace{6\baselineskip}
\blockanchor{p147-b10}
(1) (10 分) 设 $f(x, y)$ 在 $\mathbb{R}^2$ 上连续可微, 且 $\left|\frac{\partial f(x, y)}{\partial x}\right| \leq 1$, $\left|\frac{\partial f(x, y)}{\partial y}\right| \leq 1$ 对任意的 $(x, y) \in \mathbb{R}^2$ 成立. 证明: 对任意的 $(x_1, y_1), (x_2, y_2) \in \mathbb{R}^2$, 有\par
\nopagebreak[4]
% p147-b11; source PDF page 147
\blockanchor{p147-b11}
\displayformula{\left| f \left(x _ {1}, y _ {1}\right) - f \left(x _ {2}, y _ {2}\right) \right| ^ {2} \leq 2 \left(x _ {1} - x _ {2}\right) ^ {2} + 2 \left(y _ {1} - y _ {2}\right) ^ {2}.}
% p147-b12; source PDF page 147
\par\Needspace{6\baselineskip}
\blockanchor{p147-b12}
(2) (10 分) 已知函数列 $\{f_n(x, y)\}$ 在 $\mathbb{R}^2$ 上连续可微, 对任意的 $(x, y) \in \mathbb{R}^2$ 及正整数 $n$, 有\par
\nopagebreak[4]
% p147-b13; source PDF page 147
\blockanchor{p147-b13}
\displayformula{\left| \frac {\partial f _ {n} (x , y)}{\partial x} \right| \leq 1, \left| \frac {\partial f _ {n} (x , y)}{\partial y} \right| \leq 1.}
% p147-b14; source PDF page 147
\blockanchor{p147-b14}
且 $\lim_{n\to \infty}f_n(x,y) = f(x,y)$. 证明: 在 $\mathbb{R}^2$ 的任意有界闭集中, $\{f_n(x,y)\}$ 一致收敛于 $f(x,y)$.\par
